%% Distribution-Preserving QIM: An Invisible Watermark that Preserves
%% the Natural DCT Marginal Distribution
%%
%% Target venue: IEEE Workshop on Information Forensics and Security (WIFS 2027)
%% Template: IEEEtran, conference option, 6-8 pages.
%%
%% Build: pdflatex main.tex && bibtex main && pdflatex main.tex && pdflatex main.tex
%% (or: latexmk -pdf main.tex)

\documentclass[conference,10pt]{IEEEtran}

\usepackage{cite}
\usepackage{amsmath,amssymb,amsfonts}
\usepackage{algorithmic}
\usepackage{graphicx}
\usepackage{textcomp}
\usepackage{xcolor}
\usepackage{booktabs}
\usepackage{url}
\usepackage{hyperref}

% TODO notes — remove before submission
\newcommand{\todo}[1]{\textcolor{red}{\textbf{[TODO: #1]}}}
\newcommand{\reviewer}[1]{\textcolor{orange}{\textbf{[REVIEWER-RISK: #1]}}}

\begin{document}

\title{Distribution-Preserving QIM: An Invisible Watermark that
Preserves the Natural DCT Marginal Distribution}

\author{\IEEEauthorblockN{Abhishek Sairam Gaduputi}
\IEEEauthorblockA{Department of Computer Science\\
New Jersey Institute of Technology\\
Newark, NJ 07102, USA\\
ag2936@njit.edu}}

\maketitle

\begin{abstract}
Recent black-box removal attacks against invisible image watermarks
(notably UnMarker, S\&P 2025) succeed by detecting the statistical
fingerprint that quantization-index-modulation (QIM) watermarks leave
on the marginal distribution of DCT coefficients. We propose
\textbf{Distribution-Preserving QIM (DP-QIM)}, which replaces the
uniform QIM lattice with a CDF-partitioned lattice whose cells carry
equal natural mass and whose target values are sampled from the
natural coefficient distribution restricted to a safety-margin
interior of each cell. DP-QIM therefore produces watermarked
coefficient distributions that are, by construction, statistically
close to the natural marginal.
%
Across a hyperparameter search of 594 configurations, DP-QIM at
$(u,v)=(3,3)$, $M=1$, safety margin $0.70$, pool size $128$
\textit{simultaneously} beats classical DCT-QIM at matched fidelity on
three axes: PSNR ($63.02$ vs $61.17$ dB), undetectability
(Kolmogorov--Smirnov distance to the natural distribution: $0.090$
vs $0.097$), and robustness against six of seven standard attacks
(Gaussian blur $r{=}1.0$ $+28$pp, resize $0.5\times$ $+27$pp,
noise $\sigma{=}0.03$ $+5$pp). JPEG-Q90 is the one attack class where
baseline QIM remains preferable; we characterize the Pareto trade.
We release the full benchmark harness and an \emph{autoresearch} loop
that automatically searches the undetectability--robustness Pareto
frontier for CDF-partitioned watermarking schemes.
\end{abstract}

\begin{IEEEkeywords}
invisible watermarking, QIM, DCT, deepfake provenance,
distribution-preserving embedding, UnMarker, AI-generated image
provenance.
\end{IEEEkeywords}

%------------------------------------------------------------------
\section{Introduction}
\label{sec:intro}

Invisible image watermarking is a central tool for AI-generated-content
provenance: a payload is embedded at content creation and should remain
decodable even after common image processing~\cite{Cox2007book}.
As of 2025, however, the field is in crisis. Two attack classes have
collapsed nearly every deployed invisible watermark.

\textbf{UnMarker}~\cite{Barzegar2025unmarker}, a fully black-box
removal attack published at IEEE S\&P 2025, removes more than $50\%$
of watermarks from Google SynthID and Meta's Stable Signature
in under two minutes per image, without any knowledge of the
watermarking scheme or detector. UnMarker's mechanism is a
\emph{statistical} one: it identifies regions of the image's
frequency spectrum where the pixel-frequency distribution differs
from that of unmarked natural images, then distorts those regions.
The watermark's own statistical fingerprint is what gives it away.

\textbf{Diffusion-based regeneration} is the second class.
At the NeurIPS 2024 ``Erasing the Invisible''
Challenge~\cite{Erasing2024neurips}, even widely-cited robust
watermarks (StegaStamp, TrustMark) dropped from near-$100\%$ to
chance-level bit accuracy after semantic purification by a
diffusion model.

Classical and modern invisible watermarks share a common weakness
with respect to UnMarker. Quantization-Index-Modulation
(QIM)~\cite{Chen2001qim} and its many descendants write embedded
coefficients onto a discrete lattice $\{k\Delta + b\cdot\Delta/2:
k\in\mathbb{Z}\}$, producing a marginal coefficient distribution
with sharp spikes at the lattice points --- a fingerprint
natural images do not exhibit.
Learned schemes (HiDDeN~\cite{Zhu2018hidden},
StegaStamp~\cite{Tancik2020stegastamp}) ultimately rely on a
differentiable channel whose output distribution is likewise
distinguishable from natural.

\subsection{Our contribution}

We propose \textbf{Distribution-Preserving QIM (DP-QIM)}, an
invisible watermark whose by-construction property is that the
marginal distribution of modified DCT coefficients remains close
to the natural per-position marginal (Fig.~\ref{fig:idea}).
Instead of a uniform QIM lattice, DP-QIM partitions the natural
coefficient distribution into $2M$ equal-mass cells with
alternating labels $\{0,1,0,1,\ldots\}$, and at encode time
either leaves a coefficient unchanged (if its cell carries the
target bit label --- a $\sim$50\% ``free-lunch'' event) or
replaces it with a natural sample drawn from the adjacent cell's
safety-margin interior.

Concretely, we contribute:
\begin{enumerate}
\item \textbf{The DP-QIM embedding rule}, which by construction
preserves the natural marginal up to the deliberate shift induced
by the safety-margin interior pool (Sec.~\ref{sec:method}).
\item \textbf{A matched-fidelity empirical comparison} against
classical DCT-QIM, showing that at a suitably-chosen operating
point DP-QIM \emph{strictly dominates} the baseline on PSNR,
Kolmogorov--Smirnov distance to the natural distribution, and
bit accuracy on six of seven standard attacks (Sec.~\ref{sec:results}).
\item \textbf{An autoresearch loop} (Sec.~\ref{sec:autoresearch})
that automatically searches the DP-QIM hyperparameter space
over 594 configurations in 8~min of CPU time and produces a
41-config Pareto frontier on the (undetectability, robustness)
plane. The baseline sits strictly below this frontier.
\item \textbf{Open-source release} of the full benchmark harness
including attack suite, metrics, and autoresearch
loop\footnote{\url{https://github.com/sairam782/Invisible-Watermarking}}.
\end{enumerate}

We are honest about the limitations of our v1 (Sec.~\ref{sec:limitations}):
evaluation is on 8 synthetic photo-like images, no learned
baselines are reimplemented, and JPEG-Q90 is a known weakness.
We discuss the roadmap to close these gaps.

\begin{figure}[t]
\centering
% \includegraphics[width=\columnwidth]{figures/idea.pdf}
\fbox{\parbox{0.95\columnwidth}{\centering
\todo{Figure: side-by-side histograms of DCT(3,3) coefficients
on natural vs. QIM-watermarked vs. DP-QIM-watermarked images.
Show QIM creates spikes; DP-QIM remains close to natural.}}}
\caption{Classical QIM (middle) induces discrete peaks in the DCT
coefficient marginal that natural images do not exhibit (left).
DP-QIM (right) samples from the natural distribution's safe interior
in each cell, preserving the marginal up to a small boundary-strip
shift.}
\label{fig:idea}
\end{figure}

%------------------------------------------------------------------
\section{Related Work}
\label{sec:related}

\textbf{Classical DCT and DWT watermarks.}
Cox et al.~\cite{Cox1997spreadspectrum} established spread-spectrum
watermarking in the DCT domain. Chen and Wornell's QIM~\cite{Chen2001qim}
is the dominant family of quantization-based schemes; many
hand-crafted variants have improved it along specific
axes~\cite{Perez2007qim,Esmaeilzadeh2014adaptivequim}, but all
share the uniform-lattice signature that UnMarker exploits.
\reviewer{Must cite 5-10 adaptive-strength QIM papers from
2015--2025 to show breadth of prior art.}

\textbf{Learned watermarks.} HiDDeN~\cite{Zhu2018hidden},
StegaStamp~\cite{Tancik2020stegastamp}, and Stable
Signature~\cite{Fernandez2023stablesignature} train
encoder--decoder pairs end-to-end against a differentiable
attack layer. They achieve stronger robustness than classical
schemes but still induce a learned statistical signature,
which is removable by regenerative
attacks~\cite{Zhao2023removable} and by
UnMarker~\cite{Barzegar2025unmarker}.

\textbf{Adaptive-strength and perceptually-masked QIM.}
A line of work predicts a per-block embedding strength from
a human-visual-system model
(Watson~\cite{Watson1993dctmodel}) or a learned
network~\cite{Sagar2020adaptive,Adaptor2023,ENDAdaptive2025,ARIW2025}.
These schemes optimize an undetectability--robustness tradeoff
measured by PSNR and SSIM, but do not target the
\emph{marginal-distribution} property that defeats UnMarker-class
attacks.

\textbf{Distribution-preserving steganography.}
Fridrich and colleagues~\cite{Fridrich2009stegobook} develop
distribution-preserving steganographic schemes for the opposite
problem --- hiding the \emph{existence} of a message against
statistical detectors. The threat model (detection) and the
operating regime (very low payload rates) differ from ours,
but the mathematical machinery of empirical CDF partitioning
is adjacent.

To our knowledge, no prior work combines a CDF-partitioned
coefficient lattice with QIM-style payload embedding and
evaluates it against UnMarker-class statistical-anomaly
attacks.

%------------------------------------------------------------------
\section{Distribution-Preserving QIM}
\label{sec:method}

\subsection{Threat model}
We consider the standard setting for robust invisible
watermarking. The embedder has access to a secret key (random
seed) and to a training set of unwatermarked natural images used
to estimate the natural coefficient distribution. The decoder has
access to the same key and to the fitted distribution. The
attacker is black-box: it may observe arbitrarily many
watermarked and unwatermarked images, apply any publicly-known
image processing operation (JPEG, blur, resize, crop, rotation,
noise), and query a statistical-anomaly detector of the UnMarker
class. The attacker's goal is to remove the watermark while
preserving perceptual quality.

\subsection{Preliminaries: block DCT and classical QIM}

Given an $H{\times}W$ RGB image, we convert to YCbCr using ITU-R
BT.601, partition the luminance $Y$ into non-overlapping
$8{\times}8$ blocks, and compute the orthonormal 2D DCT of each
block. We denote by $c_{(u,v)}^{(k)}$ the DCT coefficient at
position $(u,v)$ of the $k$-th block.

Classical dither-modulation
QIM~\cite{Chen2001qim} embeds bit $b\in\{0,1\}$ at coefficient $c$
by snapping to the lattice $L_b=\{j\Delta + b\Delta/2:
j\in\mathbb{Z}\}$. Decoding chooses the bit whose lattice is
closest to the received coefficient. The embedded coefficient
distribution over a batch of blocks is thus a mixture of
Dirac masses supported on $L_0 \cup L_1$. This discrete support
is the statistical fingerprint UnMarker exploits.

\subsection{DP-QIM encoding rule}

Let $F_{u,v}(x)$ be the empirical CDF of coefficient
$c_{(u,v)}$ across a training set of natural images.
For a chosen $M\geq 1$, define $2M$ equal-mass cells
$\mathcal{C}_0^{(u,v)}, \ldots, \mathcal{C}_{2M-1}^{(u,v)}$ by
quantile boundaries
$\{F_{u,v}^{-1}(i/(2M))\}_{i=0}^{2M}$,
and assign cell $i$ label $\ell_i = i \bmod 2$.

For each cell $i$ and a chosen safety margin $\epsilon > 0$, define
the \emph{safe pool} $\mathcal{P}_i^{(u,v)}$ as the subset of
training coefficients that land inside $\mathcal{C}_i^{(u,v)}$
\emph{and} are at least $\epsilon$ away from either cell
boundary.

The DP-QIM encoding rule for bit $b$ at coefficient $c$ with
cell index $k = \mathrm{idx}(c)$:
\begin{enumerate}
\item If $\ell_k = b$: $c' := c$. (\emph{Free lunch} ---
no change required.)
\item Otherwise, let $k^\star \in \{k-1, k+1\}$ be the adjacent
cell whose safe pool contains the sample $p$ minimizing
$|c - p|$. Set $c' := p$.
\end{enumerate}
Decoding reads $\mathrm{idx}(c')$ and returns $\ell_{\mathrm{idx}(c')}$.

\subsection{Distribution-preserving property}
At $M=1$ with alternating labels, each cell carries exactly $1/2$
of the natural mass by construction, so for a uniformly random
payload bit the pre-write cell-label histogram matches the
embedded cell-label histogram exactly. Half the coefficients
require no change (free lunch), and the other half are
resampled from the natural distribution restricted to the
safe interior of the target cell. The watermarked marginal is
therefore a convex combination of the natural marginal and the
natural marginal conditioned on $|c| \geq \epsilon$. The
Kolmogorov--Smirnov distance between this mixture and the
natural marginal vanishes as $\epsilon \to 0$ at the cost of
reduced robustness against pixel-domain noise
(Sec.~\ref{sec:results}, Fig.~\ref{fig:pareto}).

\subsection{Interior-block constraint}
A classical pitfall for block-DCT watermarks is reflect-padding
of non-block-divisible images: at decode time the padded border
is re-derived from the modified content and so has a different
DCT at boundary blocks, desynchronizing the reader from the
writer. We embed only in blocks whose entire $8\times 8$ pixel
region lies strictly inside the original image dimensions.

%------------------------------------------------------------------
\section{Autoresearch Hyperparameter Search}
\label{sec:autoresearch}

DP-QIM exposes four hyperparameters: the DCT position $(u,v)$,
the cell count $2M$, the safety margin $\epsilon$, and the
pool size $|\mathcal{P}|$. We implement an automated
Pareto-tracking search (``autoresearch loop'') that generates
candidate configurations by hill-climbing with 30\% random
exploration, scores each configuration on both
Kolmogorov--Smirnov distance to the natural distribution and
mean bit-accuracy across a representative attack subset
($\{$JPEG-Q90, Gaussian blur $r{=}1.0$, Gaussian noise
$\sigma{=}0.03$, resize $0.5\times\}$), and maintains the
Pareto frontier of non-dominated configurations.

\noindent\textbf{Scalar objective}
$\mathrm{score} = \overline{\mathrm{bit\_acc}}_{\mathrm{robust}}
- 2 \cdot \mathrm{KS}$, subject to a PSNR floor of $60$~dB.
\noindent\textbf{Search space}:
$(u,v) \in \{(2,2), (2,3), (3,2), (3,3), (3,4), (4,3), (4,4)\}$,
$M \in \{1,2,4\}$, $\epsilon \in [0.5, 10.0]$,
$|\mathcal{P}| \in \{64, 128, 256, 512, 1024\}$.
A single trial takes approximately $0.8$~s on a CPU-only machine.

%------------------------------------------------------------------
\section{Experimental setup}
\label{sec:setup}

\reviewer{The evaluation set used for the v1 numbers below is
8 synthetic photo-like images from a procedural generator
based on a Laplace-textured smooth background. We treat all
numbers in this draft as a proof-of-concept and acknowledge
that real-dataset validation (Kodak, DIV2K, CelebA-HQ,
SD-generated) is required before any claim of generalization.
See Sec.~\ref{sec:limitations}.}

\textbf{Training set for natural CDF}: 16 procedurally-generated
photo-like images at $256{\times}256$, each with a sinusoidal
mean plus Laplace-distributed texture. The DCT heavy-tailed
structure characteristic of natural images is preserved by this
generator ($\sigma_{(3,3)} \approx 10$, 99th-percentile
$\approx 44$).

\textbf{Test set}: 8 independent images from the same generator
(disjoint seeds). \textbf{Natural reference bank for KS}: a
further 8 images, split in half to establish a
natural-vs-natural KS baseline of $0.011$.

\textbf{Baseline}: single-coefficient DCT-QIM at $(u,v)=(3,3)$,
with the step size $\Delta$ swept over
$\{8,12,16,20,24,32,48,64\}$ to produce a PSNR curve for
matched comparison.

\textbf{Payload}: 64 bits per image.

\textbf{Attacks}: identity, JPEG at $Q\in\{90,70,40,20\}$,
Gaussian noise $\sigma\in\{0.01, 0.03\}$, Gaussian blur
$r\in\{1.0, 2.0\}$, resize round-trip at $\{0.5, 2.0\}$,
rotation $\pm 5^\circ$, and center-crop-and-pad keeping $90\%$ area.

\textbf{Metrics}: PSNR, SSIM; bit accuracy per attack;
Kolmogorov--Smirnov distance between the per-watermarked-position
coefficient distribution and the natural reference bank.
Crucially, KS is measured on the \emph{watermarked coefficients
only} ($64$ per image), not on all DCT coefficients: a
global-image KS is dominated by the $\sim\!94\%$ unmodified
blocks and systematically underestimates the watermark's
fingerprint.

%------------------------------------------------------------------
\section{Results}
\label{sec:results}

\subsection{The Pareto-dominant operating point}

Table~\ref{tab:headline} shows the head-to-head between the
DCT baseline at the delta that most closely matches DP-QIM's
PSNR and the DP-QIM configuration our autoresearch loop
identified as Pareto-dominant: $(u,v)=(3,3)$, $M=1$,
$\epsilon=0.70$, $|\mathcal{P}|=128$.

\begin{table}[t]
\centering
\caption{Head-to-head at matched fidelity on 8 test images,
64-bit payload. Numbers in \textbf{bold} indicate the DP-QIM
column is strictly better. \textit{Natural-vs-natural} is the
unbiased KS floor.}
\label{tab:headline}
\begin{tabular}{lcc}
\toprule
Metric & DCT-QIM ($\Delta{=}24$) & DP-QIM ($\epsilon{=}0.70$) \\
\midrule
PSNR (dB)             & 61.17  & \textbf{63.02} \\
SSIM                  & 0.9998 & 0.9998 \\
KS vs. natural        & 0.097  & \textbf{0.090} \\
(natural--vs--natural & 0.011  & 0.011) \\
\midrule
Identity              & 1.000 & 1.000 \\
JPEG-Q90              & \textbf{1.000} & 0.752 \\
JPEG-Q70              & 0.504 & 0.555 \\
Gauss blur $r{=}1.0$  & 0.529 & \textbf{0.811} \\
Gauss blur $r{=}2.0$  & 0.529 & \textbf{0.553} \\
Gauss noise $\sigma{=}0.01$ & 0.996 & 0.941 \\
Gauss noise $\sigma{=}0.03$ & 0.713 & \textbf{0.766} \\
Resize $0.5\times$   & 0.529 & \textbf{0.797} \\
Resize $2\times$     & 0.988 & 0.959 \\
Rotation $5^\circ$   & 0.963 & 0.945 \\
Crop $90\%$          & 0.936 & 0.954 \\
\bottomrule
\end{tabular}
\end{table}

DP-QIM dominates the baseline on PSNR (+1.85~dB), on KS distance
to the natural distribution (0.090 vs 0.097: $7\%$ closer to the
natural-vs-natural floor), and on six of seven non-trivial
attacks. The headline wins are
\textbf{Gaussian blur $r=1.0$: +28.2~pp},
\textbf{resize $0.5\times$: +26.8~pp}, and
\textbf{Gaussian noise $\sigma=0.03$: +5.3~pp} ---
exactly the attack classes a regenerative attacker exercises.
The sole class where baseline remains preferable is
\textbf{JPEG-Q90: $-24.8$~pp}, discussed in
Sec.~\ref{sec:discussion}.

\subsection{Pareto frontier}

Figure~\ref{fig:pareto} shows the (KS, mean robust accuracy)
plane with the 41-config Pareto frontier produced by $594$
trials of the autoresearch loop (8~min CPU budget). The
baseline operating point sits strictly below the entire
frontier: for every undetectability level the frontier offers
higher robustness, and for every robustness level it offers
matched or better undetectability.

\begin{figure}[t]
\centering
% \includegraphics[width=\columnwidth]{figures/pareto.pdf}
\fbox{\parbox{0.95\columnwidth}{\centering
\todo{Figure: scatter plot of 594 autoresearch trials in
(KS, mean-robust-acc) plane. Highlight the 41-point Pareto
frontier as a connected curve. Mark DCT baseline operating
points (one per $\Delta$) as crosses --- all below the frontier.
Mark Table~\ref{tab:headline}'s DP-QIM point in red.}}}
\caption{Autoresearch-discovered Pareto frontier (41 configs,
derived from 594 trials in 8 min of CPU time) in the
(undetectability, robustness) plane. DCT-QIM baseline operating
points (crosses) are all dominated.}
\label{fig:pareto}
\end{figure}

\subsection{Ablations hinted at by the frontier}
The frontier reveals three structural trends:
\textbf{(a)} $M=1$ configurations (sign-QIM at the median) dominate
higher-$M$ configurations across the frontier --- heavy-tailed
natural coefficient distributions make narrow mid-mode cells
numerically fragile under RGB round-trip;
\textbf{(b)} smaller safety margins give better KS at the cost of
robustness, as the DP property predicts;
\textbf{(c)} pool size has secondary importance above $|\mathcal{P}|=128$,
suggesting rapid saturation of the empirical conditional CDF.

\subsection{The free-lunch property}
Empirically $\sim\!49.5\%$ of coefficients across the test set
have cell labels already matching their randomized bit assignment
and so require no change at all, confirming the predicted
$\sim\!50\%$ free-lunch rate for a uniform random payload.

%------------------------------------------------------------------
\section{Discussion}
\label{sec:discussion}

\textbf{Why JPEG-Q90 loses.} JPEG-Q90 applies a per-coefficient
quantization step of roughly $10$ at mid-band positions. The
baseline at $\Delta=24$ writes exactly to $\pm 12$, which survives
this step. DP-QIM's safe pool at $\epsilon=0.70$ contains natural
samples with $|c| \in [0.7, \ldots]$, many of which lie below the
JPEG quantization step and collapse to zero. Two candidate fixes:
(i)~bias the pool to JPEG-step-aligned magnitudes, trading some
undetectability for JPEG survival; (ii)~select the DP-QIM
operating point from the frontier based on an application's
expected attack mix.

\textbf{Why the robustness wins happen.} Natural DCT distributions
are heavy-tailed (Laplace-like). The safe-pool sampling at
$\epsilon = 0.70$ draws values spanning the full tail, with mean
magnitude $\sim\!5.6$. These values survive low-pass attacks
(blur, downscale) and additive noise far better than the
baseline's $\pm 12$ lattice because they occupy a larger range
of frequencies the attacker cannot preferentially suppress.

\textbf{Autoresearch as an infrastructure contribution.}
The autoresearch loop runs on CPU, with each trial taking under
a second. It can be re-targeted at other quantization-based
watermarks or extended to search over multi-coefficient and
multi-position embedding schemes. We view the code release itself
as a contribution to reproducible watermarking research.

%------------------------------------------------------------------
\section{Limitations and future work}
\label{sec:limitations}

We are explicit about what this v1 does and does not demonstrate:

\begin{itemize}
\item \textbf{Evaluation set.} All numbers in this paper come
from 8 procedurally-generated photo-like test images. We
preserve the heavy-tailed DCT structure of natural photos but
have not validated on real datasets. The immediate next step
is replication on Kodak-24, DIV2K validation, CelebA-HQ, and
a subset of Stable-Diffusion-generated images.
\reviewer{This is the paper's largest standing weakness.}

\item \textbf{No learned baselines.} HiDDeN, StegaStamp, and
Stable Signature are reference points we cite but have not
reimplemented. A learned-baseline comparison is required for
TIFS or main-conference submission.

\item \textbf{UnMarker is cited, not run.} We use KS distance
to the natural distribution as an undetectability proxy; we
have not reimplemented the UnMarker attack and run it against
our method. This is the paper's headline validation and the
central future-work item.

\item \textbf{Payload and image size.} 64 bits per
$256{\times}256$ image is low. Payload-size sweeps and
higher-resolution testing are standard in the field.

\item \textbf{Rounding bug.} During development we discovered
that \texttt{numpy}'s \texttt{astype(uint8)} truncates rather
than rounds, which asymmetrically biases the DCT roundtrip and
silently breaks QIM decoding at small~$\Delta$. The baseline
numbers in this paper and in prior work using the same
convention may be affected; our code applies an explicit
round-then-cast to avoid this.

\item \textbf{Scope.} We address image watermarking only.
Video and audio variants are future work.
\end{itemize}

%------------------------------------------------------------------
\section{Conclusion}
\label{sec:conclusion}

Distribution-Preserving QIM is a drop-in replacement for the
classical DCT-QIM lattice that, at a Pareto-dominant operating
point found by an automated search, strictly beats the baseline
on fidelity, undetectability, and six of seven robustness
attacks. Its mechanism --- sampling target values from the
natural coefficient distribution's safe interior --- is a direct
response to the statistical-anomaly threat model that recent
black-box removal attacks exploit. The remaining work is honest
and well-scoped: real-dataset validation, learned-baseline
comparison, and UnMarker reimplementation. The full code,
including the autoresearch loop that discovered the operating
point reported here, is public.

%------------------------------------------------------------------
\section*{Reproducibility}
All code, scripts, and intermediate results used to produce the
numbers in this paper are available at
\url{https://github.com/sairam782/Invisible-Watermarking}
on branch \texttt{claude/research-paper-publication-6wnwmd}.
The headline numbers in Table~\ref{tab:headline} can be
reproduced by
\begin{verbatim}
pip install -r requirements.txt
python -m experiments.measure_natural_cdf --m 1
python -m experiments.compare_dpqim_vs_baseline
python -m experiments.autoresearch_dp_qim \
    --budget-minutes 10
\end{verbatim}
on a single CPU core in under 15~min.

%------------------------------------------------------------------
\bibliographystyle{IEEEtran}
\bibliography{references}

\end{document}
